\section{The actual arithmetic local charts and their residual representations}
\label{sec:arithmetic-local-charts-v055}

\paragraph{What the local construction must preserve.}
The component argument uses the entry algebra of the supplied universal
representation, not an arbitrary ring with a similar spectrum.  At a
closed point it passes to a specified coefficient extension and a
complete local ring.  We prove that the original universal property
makes this ring the actual completion, prove its flatness, and identify
the representation and all three entry ideals after that map.  We also
construct these local rings at every original permitted closed point.
The residual argument treats ranks two and three separately where
their matrix patterns differ.

The upstream results RTC7 and RTC9 are available at their full original
scope.  Thus the actual weight-zero coefficient ring is complete
Noetherian local, the actual graded entry algebra is Noetherian and
finite type over it, and all entries and ordinary characters have the
required adic continuity.  None of those results assumes a later
\(\mathrm{FinA}\), finite-point, propagation or automorphic conclusion.

\paragraph{Notation.}
In the arithmetic statements let \(O\) be the original complete
discrete valuation ring, with finite residue field \(k\) of
characteristic \(\ell\), and let \(F\) satisfy the original
\(F.\mathrm{Ok}(C_0)\).  Fix \(n=2\) or \(3\) and a supplied
\(S:\mathrm{Space}(O,F,n)\).  Write \(\mathcal R=S.R\),
\(B=S.B\), \(A=S.A\), and \(u=S.u\).  Their inclusions are
the original inclusions in \(\mathcal R\).  Put
\[
 \mathfrak m_B=B\cap\mathfrak m_{\mathcal R},\qquad
 x_0=\mathfrak m_B A,\qquad
 \pi(\mathfrak p)=\mathfrak p\cap B.
                                                        \tag{ALC.1}
\]
The weights are \(w(2)=(1,-1)\) and \(w(3)=(1,0,-1)\).
The ideals \(\mathfrak t_+,\mathfrak t_-\) are generated by the
entries below and above the diagonal.  The ideal \(\mathfrak e\)
is generated by entries whose weight difference is odd.  Thus
\[
 \mathrm{TI}(\mathfrak p)
 \Longleftrightarrow
 \mathfrak t_+\not\subseteq\mathfrak p,
 \quad\mathfrak t_-\not\subseteq\mathfrak p,
 \quad(n=3\Rightarrow\mathfrak e\not\subseteq\mathfrak p).
                                                        \tag{ALC.2}
\]
An original \(L:S.\mathrm{Loc}(x)\) contains a finite coefficient
extension \(O\to O'\), with
\(\mathfrak m_OO'=\mathfrak m_{O'}\), a coefficient algebra
\(T=L.R\) over \(O'\), and the specified \(O\)-algebra map
\(\iota:A\to T\).  It requires
\(\iota^{-1}\mathfrak m_T=x\) and the residue-compatible
universal property used explicitly below.  Its representation is
\(\rho_L=u\otimes\iota\); for an ideal \(P\subset T\) write
\(L.\mathrm{im}(P)=\iota^{-1}P\) and
\(L.\pi(P)=\pi(\iota^{-1}P)\).

The first completion statements have the wider original algebraic
scope: \(O\) is merely a commutative local ring.  A coefficient
algebra over it means exactly \(\mathrm{IsCO}\): a complete
Noetherian local \(O\)-algebra, with a local structure map and
the specified residue identification with \(\operatorname{Res}(O)\).
No finiteness of that residue field is needed in these statements.

\subsection{The universal property determines every finite quotient}

% BEGIN ALC1
\paragraph{ALC1. Locality of coefficient maps and the nilpotent coefficient quotients.}
\label{claim:ALC1}
Every \(O\)-algebra map between coefficient algebras over a
commutative local ring \(O\) is local, preserves the specified
residue maps, and is continuous for all powers of the maximal
ideals.  If \(D\) is a Noetherian \(O\)-algebra and
\(\varphi:D\to k=\operatorname{Res}(O)\) is an
\(O\)-algebra map, put \(y=\ker\varphi\).  Then
\(D/y^N\) is a coefficient algebra for every \(N>0\), with
nilpotent maximal ideal \(y/y^N\) and residue map induced by
\(\varphi\).

\begin{proof}
Let \(h:R\to C\) be such a map and \(r\in\mathfrak m_R\).
Choose \(a\in O\) whose image in \(C\) has the same residue
as \(h(r)\).  If \(a\) were a unit, then \(r-a\) would be
a unit in \(R\), whereas \(h(r-a)\in\mathfrak m_C\).
This is impossible.  Thus \(a\in\mathfrak m_O\), and the
local structure map gives \(h(r)\in\mathfrak m_C\).
For arbitrary \(r\), choose \(a\) with \(r-a\in\mathfrak m_R\)
and apply this result.  It proves
\(\operatorname{res}_C(h(r))=\operatorname{res}_R(r)\).
Also \(h(\mathfrak m_R^N)\subseteq\mathfrak m_C^N\) for
every \(N\), which is precisely the required continuity.

The map \(\varphi\) is surjective because its restriction to
\(O\) is the residue map.  Hence \(y\) is maximal.  In
\(D/y^N\), the ideal \(y/y^N\) is nilpotent and its quotient
is \(k\).  An element outside this ideal has an inverse modulo
the ideal; the finite geometric series corrects that inverse to
an actual inverse.  Consequently this is the unique maximal ideal.
Its \(N\)th power is zero, so its adic completion is the ring
itself and its adic filtration is separated.  Noetherianity follows
from that of \(D\).  The map from \(O\) induces the original
residue map and sends nonunits into the maximal ideal.  These are
all the \(\mathrm{IsCO}\) clauses.  In particular this applies
to \(R/\mathfrak m_R^N\) whenever \(R\) is a coefficient
algebra.  It also applies with
\(\varphi=\operatorname{res}_R\circ\iota\) for a supplied
map \(\iota:D\to R\).
\end{proof}
% END ALC1

% BEGIN ALC2
\paragraph{ALC2. The two quotient maps are actual inverses.}
\label{claim:ALC2}
Let \(D\) be a Noetherian \(O\)-algebra, let \(R\) be a
coefficient algebra, and let \(\iota:D\to R\) be an
\(O\)-algebra map.  Assume the original universal property:
for every coefficient algebra \(C\) and every \(f:D\to C\)
with \(\operatorname{res}_C f=\operatorname{res}_R\iota\),
there is a unique \(O\)-algebra map \(g:R\to C\) satisfying
\(g\iota=f\).  Put \(y=\iota^{-1}\mathfrak m_R\).
For every \(N>0\), the map induced by \(\iota\) is an
isomorphism
\[
 \alpha_N:D/y^N\xrightarrow{\ \sim\ }R/\mathfrak m_R^N.
                                                        \tag{ALC.3}
\]
These isomorphisms commute with every reduction to a smaller
positive exponent.

\begin{proof}
ALC1 makes \(D/y^N\) a coefficient algebra with the required
residue map.  Apply the universal property to its quotient map
to obtain \(g_N:R\to D/y^N\), with
\(g_N\iota=q_D\).  By ALC1 the map \(g_N\) is local.
The target maximal ideal has \(N\)th power zero, so
\(g_N(\mathfrak m_R^N)=0\).  It therefore factors as
\(\beta_N:R/\mathfrak m_R^N\to D/y^N\).
Conversely \(\iota(y^N)\subseteq\mathfrak m_R^N\) defines
\(\alpha_N\).

The equality \(\beta_N\alpha_N=1\) follows by evaluating on
the surjective image of \(D\): both composites send \(d\)
to its class modulo \(y^N\).  To prove the other equality,
use the universal property with the coefficient algebra
\(R/\mathfrak m_R^N\).  The two maps
\(q_R\) and \(\alpha_Ng_N\) from \(R\) agree after
\(\iota\), and their common prescribed map has the required
residue.  Uniqueness gives \(\alpha_Ng_N=q_R\), hence
\(\alpha_N\beta_N=1\).  Reduction compatibility for
\(\alpha_N\) follows directly from their definition on \(D\);
their inverses consequently have the same compatibility.
\end{proof}
% END ALC2

% BEGIN ALC3
\paragraph{ALC3. Full completion from the original universal property.}
\label{claim:ALC3}
Under exactly the hypotheses of ALC2, \(\iota\) satisfies the
original \(\mathrm{IsCplAt}(y,R)\) record:
\[
 \iota^{-1}(\mathfrak m_R^N)=y^N,
 \qquad
 \forall r\in R\ \exists d\in D:
           r-\iota(d)\in\mathfrak m_R^N
                     \quad(N\ge0),                  \tag{ALC.4}
\]
and \(R\) is Noetherian and complete for its maximal ideal.
This proves the full original \(\mathrm{cpl\_of\_univ}\).

\begin{proof}
For \(N>0\), injectivity of (ALC.3) is exactly the first
identity in (ALC.4), and surjectivity is exactly the second.
For \(N=0\), the first identity says that the inverse image
of the unit ideal is the unit ideal; the second holds with
\(d=0\).  Noetherianity and completeness are already fields
of the supplied coefficient-algebra structure.  These are
literally the four fields in \(\mathrm{IsCplAt}\).

In particular, the compatible isomorphisms (ALC.3), followed
by completeness and separatedness of \(R\), identify it with
\(\varprojlim_N D/y^N\).  This identification sends the
original element \(d\) to \(\iota(d)\); it does not replace
the original map with an unmarked isomorphism of rings.
\end{proof}
% END ALC3

% BEGIN ALC4
\paragraph{ALC4. The full original arithmetic local-completion theorem.}
\label{claim:ALC4}
Return to the arithmetic data and let \(L:S.\mathrm{Loc}(x)\)
be any supplied local object.  For every \(O'\)-algebra map
\[
 \iota':D=O'\otimes_OA\longrightarrow T,
 \qquad \iota'(1\otimes a)=\iota(a)\quad(a\in A),
                                                        \tag{ALC.5}
\]
the original \(\mathrm{IsCplAt}\) record holds at
\(y=(\iota')^{-1}\mathfrak m_T\).  This is the full
original \(\mathrm{loc\_cpl}\), for both ranks and every
original \(x,L,\iota'\).

\begin{proof}
By RTC7, \(A\) is Noetherian.  The specified \(O'\) is a
finite \(O\)-module, so \(D\) is a finite \(A\)-module:
tensor a finite set of \(O\)-module generators by \(A\).
It is therefore a Noetherian ring.

Let \(C\) be any coefficient algebra over \(O'\), and let
\(f:D\to C\) be an \(O'\)-algebra map whose residue agrees
with \(\iota'\).  Give \(C\) its composite \(O\)-algebra
structure and the resulting scalar tower.  Restrict \(f\)
along \(a\mapsto1\otimes a\).  By (ALC.5), its residue
agrees with \(\iota:A\to T\).  The actual \(L.\mathrm{univ}\)
therefore supplies a unique \(O'\)-algebra map \(g:T\to C\)
agreeing on \(A\).

For every pure tensor \(b\otimes a\), \(O'\)-linearity
and the equality on \(A\) give
\(g(\iota'(b\otimes a))=b f(1\otimes a)=f(b\otimes a)\).
Pure tensors additively generate \(D\), so \(g\iota'=f\).
Any other such map agrees on the original \(A\), hence is
equal by \(L.\mathrm{univ}\).  We have verified every
antecedent of ALC3 with coefficient ring \(O'\), the given
\(T\), and this very \(\iota'\).  Its conclusion is the
required record.
\end{proof}
% END ALC4

\subsection{Completion is flat, and the primes remain attached to the original map}

% BEGIN ALC5
\paragraph{ALC5. Constructing the Noetherian completion and its exact powers.}
\label{claim:ALC5}
Let \(E\) be a Noetherian commutative ring, let
\(I=(a_1,\ldots,a_r)\), and define
\(\widehat E=\varprojlim_N E/I^N\).  Then
\[
 \ker(\widehat E\to E/I^N)=I^N\widehat E,
 \qquad \widehat E/I^N\widehat E=E/I^N.          \tag{ALC.6}
\]
The ring \(\widehat E\) is Noetherian, complete and separated
for \(I\widehat E\), and \(E\) is dense in it.  If \(I\)
is maximal, \(\widehat E\) is local with maximal ideal
\(I\widehat E\), and the source's particular
\(\mathrm{Cpl}(I)=\widehat{E_I}\) has the full
\(\mathrm{IsCplAt}(I,\mathrm{Cpl}(I))\) record.  The last
assertion is the original \(\mathrm{cpl\_is}\).

\begin{proof}
Formal substitution defines a ring map
\[
 E[[X_1,\ldots,X_r]]\longrightarrow\widehat E,
                \qquad X_i\longmapsto a_i.           \tag{ALC.7}
\]
It is defined by finite truncations: terms of total degree at
least \(N\) vanish modulo \(I^N\).  Thus multiplication
and addition are legitimate in every finite quotient.

For \(z\in\widehat E\), choose its constant coefficient
in \(E\) to match \(z\) modulo \(I\).  Suppose terms of
degrees below \(N\) have been chosen.  Their discrepancy
with \(z\), reduced modulo \(I^{N+1}\), belongs to
\(I^N/I^{N+1}\).  The degree-\(N\) monomials in the
\(a_i\) generate that module over \(E\), so choose their
coefficients to correct the discrepancy.  Continuing gives a
formal series mapping to \(z\).  This proves surjectivity of
(ALC.7).  If \(z\) is zero modulo \(I^N\), start with all
terms of degree below \(N\) equal to zero.  The resulting
series lies in \((X_1,\ldots,X_r)^N\), and hence
\(z\in I^N\widehat E\).  The reverse inclusion is immediate.
This proves (ALC.6), including the case \(N=0\).

The finite-variable formal power-series ring over a Noetherian
ring is Noetherian; its surjective image (ALC.7) is consequently
Noetherian.  The exact kernel identities identify the adic
quotients of \(\widehat E\) with the inverse system defining
it.  They prove completeness and separatedness, and their
surjectivity from \(E\) proves density.

When \(I\) is maximal, the quotient by \(I\widehat E\) is
the field \(E/I\).  An element invertible in that quotient
has a lift of its inverse whose product differs from one by
an element of \(I\widehat E\).  The convergent geometric
series corrects the lift.  Therefore every element outside
\(I\widehat E\) is a unit, proving the local assertion.

For the source's particular completion, observe that
\(E/I^N\) is already local when \(I\) is maximal and
\(N>0\): its maximal ideal is nilpotent and its quotient
is a field.  Every element of \(E\setminus I\) therefore
becomes a unit in it.  Localization consequently gives
\[
        E/I^N\simeq E_I/(IE_I)^N.
\]
These maps respect all reductions and the original maps
from \(E\).  Passing to the inverse limits identifies
\(\widehat E\) with \(\widehat{E_I}\) as an
\(E\)-algebra.  The exact quotient identities and density
above give all four \(\mathrm{IsCplAt}\) fields for this
specific ring and map.
\end{proof}
% END ALC5

% BEGIN ALC6
\paragraph{ALC6. The induced filtration on a finite submodule.}
\label{claim:ALC6}
For a finite module \(M\) over a Noetherian ring \(E\), a
submodule \(N\subset M\), and an ideal \(I\subset E\),
there is \(c\ge0\) such that
\[
 N\cap I^qM=I^{q-c}(N\cap I^cM)\quad(q\ge c).
                                                        \tag{ALC.8}
\]
In particular the filtrations \(I^qN\) and \(N\cap I^qM\)
are cofinal.

\begin{proof}
The Rees algebra \(\bigoplus_{q\ge0}I^qt^q\) is generated
over \(E\) by finitely many elements \(a_it\), hence is
Noetherian.  The Rees module \(\bigoplus_{q\ge0}I^qMt^q\)
is finite over it, generated by a finite set of generators of
\(M\) in degree zero.  Its graded submodule
\(\bigoplus_{q\ge0}(N\cap I^qM)t^q\) therefore has a
finite homogeneous generating family.  Choose \(c\) at
least their degrees.

A generator of degree \(d\le c\) contributes in degree
\(q\ge c\) an element of
\(I^{q-d}(N\cap I^dM)\).  This lies in
\(I^{q-c}(N\cap I^cM)\).  The opposite inclusion follows
from multiplication inside \(M\), giving (ALC.8).
Now \(I^qN\subset N\cap I^qM\), while
\(N\cap I^{q+c}M\subset I^qN\).  These are the two
cofinality inclusions.  No assumption that \(I\) lies in a
Jacobson radical was used.
\end{proof}
% END ALC6

% BEGIN ALC7
\paragraph{ALC7. Exactness of finite-module completion and its tensor description.}
\label{claim:ALC7}
Completion for an ideal \(I\) of a Noetherian ring \(E\)
is exact on short exact sequences of finite modules.  For
every finite \(E\)-module \(M\), its canonical map gives
\[
               \widehat E\otimes_E M\simeq\widehat M.
                                                        \tag{ALC.9}
\]

\begin{proof}
Consider \(0\to N\to M\to Q\to0\).  ALC6 identifies
the completion of \(N\) in its own filtration with that in
the filtration induced from \(M\), so its map into
\(\widehat M\) is injective.

For surjectivity onto \(\widehat Q\), choose successive
representatives \(q_j\in Q\) with
\(q_{j+1}-q_j\in I^jQ\).  Having lifted \(q_j\) to
\(m_j\), lift the difference to an element of \(I^jM\)
and add it to \(m_j\).  This produces a coherent sequence
\(m_j\) lifting the prescribed sequence.  Thus
\(\widehat M\to\widehat Q\) is surjective.

If \(z\in\widehat M\) maps to zero, represent it by
\(m_j\in M\) modulo \(I^jM\).  Its image in \(Q\)
lies in \(I^jQ\); subtract a lift from \(I^jM\) to obtain
\(n_j\in N\).  The differences of these representatives
belong to \(N\cap I^jM\).  Cofinality in ALC6 makes them
a coherent element of \(\widehat N\), with image \(z\).
This proves exactness also at the middle term.

Choose a finite presentation \(E^s\to E^t\to M\to0\),
which exists by Noetherianity.  Applying the just proved
exactness, and using that completion commutes with finite
free modules, identifies \(\widehat M\) with the cokernel
of the same matrix over \(\widehat E\).  Right exactness
of tensor product identifies \(\widehat E\otimes_E M\)
with that cokernel as well.  These identifications agree on
\(1\otimes m\), and hence give the canonical (ALC.9).
\end{proof}
% END ALC7

% BEGIN ALC8
\paragraph{ALC8. Flatness of completion for the actual algebra structure.}
\label{claim:ALC8}
The completion \(\widehat E\) of every Noetherian ring at
every ideal is flat over \(E\).  More generally, if a
Noetherian ring \(D\), an ideal \(y\), and a local
\(D\)-algebra \(R\) satisfy the original
\(\mathrm{IsCplAt}(y,R)\), then this very \(D\)-algebra
\(R\) is flat.

\begin{proof}
For every finitely generated ideal \(J\subset E\), complete
the injection \(J\hookrightarrow E\).  ALC7 and its
canonical tensor description give an injection
\[
              J\otimes_E\widehat E\longrightarrow\widehat E.
                                                        \tag{ALC.10}
\]
The ideal criterion for flatness states that an \(E\)-module
is flat precisely when these maps are injective for all
finitely generated ideals.  It is the finite-relation form
of tensor exactness: relations between a finite set of
coefficients can be resolved in their generated ideal, and
then successively in a finite free presentation.  Applying
that criterion to (ALC.10) proves the first assertion.

For the second assertion, the exact power-contraction and
density fields of \(\mathrm{IsCplAt}\) give canonical
isomorphisms \(D/y^N\simeq R/\mathfrak m_R^N\) for
every positive \(N\).  They commute with reduction.  The
complete separated ring \(R\) is therefore isomorphic to
\(\varprojlim_ND/y^N\) as a \(D\)-algebra, with its
original map from \(D\).  The first assertion applies to
that completion.  Transport of modules through this
\(D\)-algebra isomorphism proves flatness for the specified
structure on \(R\).
\end{proof}
% END ALC8

% BEGIN ALC9
\paragraph{ALC9. Full prime and minimal-prime return for a completed local chart.}
\label{claim:ALC9}
Let \(D\) be Noetherian, \(y\) maximal, and \(R\) a
local \(D\)-algebra satisfying \(\mathrm{IsCplAt}(y,R)\).
Besides flatness, the following all hold:
\begin{enumerate}
\item every prime \(p\subseteq y\) of \(D\) is the
contraction of a prime of \(R\);
\item every prime of \(R\) contracts into \(y\);
\item every minimal prime of \(R\) contracts to a minimal
prime of \(D\);
\item for every ideal \(J\subset D\), every prime minimal
over \(JR\) contracts to a prime minimal over \(J\).
\end{enumerate}
Together these prove the full original theorem \(\mathrm{c3a}\).

\begin{proof}
ALC8 proves flatness, and the power-contraction identity with
exponent one gives \(\mathfrak m_R\cap D=y\).
We spell out the prime argument so that it applies to the
actual map and to every ideal \(J\).

For a prime \(Q\subset R\), put \(p=Q\cap D\).
The induced map \(D_p\to R_Q\) is flat and local.  It is
faithfully flat: a proper ideal of \(D_p\) extends into the
maximal ideal of \(R_Q\), and hence remains proper.  Any
nonzero module has a nonzero cyclic submodule; flatness
preserves its injection and the preceding proper-ideal
assertion keeps that cyclic module nonzero after tensoring.
If \(p_0\subseteq p\), the fiber of \(R_Q\) over the
prime \(p_0D_p\) is consequently a nonzero ring.  Choose
a prime of that fiber.  Its contraction to \(R\) gives
\(Q_0\subseteq Q\) with \(Q_0\cap D=p_0\).
This proves the needed going-down assertion directly.

Apply it first to \(Q=\mathfrak m_R\) to obtain the first
claim.  The second follows from \(Q\subseteq\mathfrak m_R\).
For the fourth, let \(Q\) be minimal over \(JR\), and
let \(p_0\) be a prime containing \(J\) and contained in
\(Q\cap D\).  Going down produces \(Q_0\subseteq Q\)
over \(p_0\).  Since \(J\subseteq p_0\), this prime
contains \(JR\).  Minimality forces \(Q_0=Q\), hence
\(p_0=Q\cap D\).  This is exactly minimality over \(J\).
The third claim is its case \(J=0\).  No choice of a
different completion or of a different prime above \(y\)
has entered the argument.
\end{proof}
% END ALC9

% BEGIN ALC10
\paragraph{ALC10. Finite-residue contractions and the original coefficient extension.}
\label{claim:ALC10}
If a ring \(D\) maps to a local ring with finite residue
field, the inverse image of its maximal ideal is maximal.
If \(O,O'\) are discrete valuation domains and
\(\mathfrak m_OO'=\mathfrak m_{O'}\), then the specified
map \(O\to O'\) is injective and \(O'\) is flat over \(O\).

\begin{proof}
For the first assertion the image of \(D\) in the finite
residue field is a finite integral domain.  Multiplication
by any nonzero element of that image is injective, hence
surjective, so the image is a field.  Its kernel is the
stated contraction, proving maximality.

For the second choose a uniformizer \(\varpi\) of \(O\).
Its image generates \(\mathfrak m_{O'}\), so is nonzero:
otherwise the discrete valuation ring \(O'\) would be a
field.  A nonzero element of \(O\) is a unit times a power
of \(\varpi\).  Its image is nonzero in the domain \(O'\),
proving injectivity.  Thus \(O'\) is torsion-free over \(O\).
Every nonzero ideal of a DVR is principal, and its tensor
map into \(O'\) is multiplication by a nonzero element.
It is injective by torsion-freeness.  The same ideal
criterion used in ALC8 proves flatness.
\end{proof}
% END ALC10

% BEGIN ALC11
\paragraph{ALC11. Full original flatness and maximality of the arithmetic chart.}
\label{claim:ALC11}
For every original \(x\) and every supplied
\(L:S.\mathrm{Loc}(x)\), the map \(\iota:A\to T\) is
flat, \(x\) is maximal, and
\[
                    \mathfrak m_B\subseteq\pi(x).
                                                        \tag{ALC.11}
\]
In fact equality holds in (ALC.11).  This supplies every
field of the full original \(\mathrm{loc\_flat}\) for
both allowed ranks.

\begin{proof}
Use the canonical \(O'\)-algebra map
\(\iota':O'\otimes_OA\to T\), sending
\(b\otimes a\) to \(b\iota(a)\).  ALC4 identifies it
as the original completion map, and ALC8 proves it flat.
ALC10 and the supplied equality
\(\mathfrak m_OO'=\mathfrak m_{O'}\) show that \(O'\)
is flat over \(O\).  Tensoring therefore makes
\(A\to O'\otimes_OA\), \(a\mapsto1\otimes a\),
flat.  The composite is exactly \(\iota\), so it is flat.

The finite residue field of \(T\) is the specified residue
field of \(O'\).  Applying the first part of ALC10 to
\(\iota\) shows that its maximal-ideal contraction \(x\)
is maximal.  Apply the same assertion to the composite
\(B\to A\to T\).  Its contraction \(\pi(x)\) is a
maximal ideal of the actual local ring \(B\).  RTC7 says
that this ring's maximal ideal is exactly \(\mathfrak m_B\).
Thus \(\pi(x)=\mathfrak m_B\), proving (ALC.11).
\end{proof}
% END ALC11

\subsection{The residual types in ranks two and three}

% BEGIN ALC12
\paragraph{ALC12. Opposite triangular directions cannot both survive on the special fiber.}
\label{claim:ALC12}
Let \(K\) be any field and \(f:A\to K\) any ring map,
with no prescribed coefficient-field structure on \(K\).
If \(\mathfrak m_B\subseteq\pi(\ker f)\), then
\[
             \mathfrak t_+\subseteq\ker f
             \quad\hbox{or}\quad
             \mathfrak t_-\subseteq\ker f.
                                                        \tag{ALC.12}
\]
This holds for both original ranks, including rank three.

\begin{proof}
Suppose neither inclusion holds.  Since these are generated
entry ideals, choose a below-diagonal entry \(a=u(g)_{ij}\)
and an above-diagonal entry \(b=u(h)_{rs}\) with
\(f(a)\ne0\) and \(f(b)\ne0\).  In both allowed ranks
the weights strictly decrease with the index.  Thus the
weight of \(a\) is \(-p\), and that of \(b\) is \(q\),
for positive integers \(p,q\).

The product \(a^qb^p\) has weight zero.  By the definition
of the actual coefficient algebra it belongs to \(B\).
Both entries are off the diagonal, so their images in the
residue field of the original \(\mathcal R\) are zero.
Consequently \(a^qb^p\in\mathfrak m_B\).  The assumed
contraction condition gives
\[
                       f(a)^q f(b)^p=0.
\]
This is impossible in a field when both factors are nonzero.
The argument uses an actual weight-zero entry monomial; it
does not use the rank-two triangular point realization, nor
does it require a finite or algebraic residue-field target.
\end{proof}
% END ALC12

% BEGIN ALC13
\paragraph{ALC13. The marked difference gives every coordinate projector.}
\label{claim:ALC13}
The actual \(B\)-linear span of \(u(G_0)\) contains each
coordinate projector \(E_{ii}\).  More generally the full
original \(\mathrm{idem}\) statement holds: for any original
coefficient algebra \(C\), any commutative \(C\)-algebra
\(K\), and any normalized \(\rho\) with
\(\mathrm{Integral}(O,C,\rho)\), the \(C\)-span of
\(\rho(G_0)\) contains all \(E_{ii}\).

\begin{proof}
Here is the elementary matrix calculation behind both assertions.
The linear span of a multiplicative matrix representation is
an algebra: products of two generators are again generators,
and its identity is the matrix at the group identity.
Suppose \(2\) and \(\zeta\) are units in the scalar ring
and that this algebra contains
\[
 \rho(g_0)-\rho(g_0^*)
             =\zeta\operatorname{diag}(w(n)_i).
                                                        \tag{ALC.13}
\]
Multiplying by \(\zeta^{-1}\) puts
\(Y=\operatorname{diag}(w(n)_i)\) in the algebra.  For
rank two its projectors are
\[
                  E_{00}=(1+Y)/2,
                  \qquad E_{11}=(1-Y)/2.
                                                        \tag{ALC.14}
\]
For rank three they are
\[
 E_{00}=(Y^2+Y)/2,
 \qquad E_{11}=1-Y^2,
 \qquad E_{22}=(Y^2-Y)/2.
                                                        \tag{ALC.15}
\]
Each formula is checked on the diagonal entries \(1,-1\)
or \(1,0,-1\); off-diagonal entries vanish.  In particular
no division by \(3\) occurs.

For the original \(\mathrm{idem}\) statement, the
\(\mathrm{Integral.ze}\) field supplies \(\zeta\in C\)
whose residue is \(\Psi(g_0)-\Psi(g_0)^{-1}\), together
with the diagonal entries of (ALC.13).  The off-diagonal
entries vanish by the normalized \(\mathrm{diag}\) condition.
The residue of \(\zeta\) is nonzero because
\(\Psi(g_0)^2\ne1\), the original marked-element field
of \(F.\mathrm{Ok}\).  Thus \(\zeta\) is a unit in
the local ring \(C\).  Since \(\ell\ne2\), the residue
of \(2\) is also nonzero, so \(2\) is a unit.  Applying
(ALC.14) or (ALC.15) proves the full claimed span statement
over every original \(C,K\).  The additional trace and
ordinary-character fields of \(\mathrm{Integral}\) need
no strengthening.

For the actual representation over \(A\), a full
\(\mathrm{Integral}\) record need not be imported.  Set
\(\zeta=u(g_0)_{00}-u(g_0)_{n-1,n-1}\in B\).
Its residue is the same nonzero difference.  RTC7 makes
\(B\) a local coefficient ring, so \(2,\zeta\) are units
there.  Polarization gives
\(u(g_0^*)_{ij}=u(g_0)_{j^{\rm rev},i^{\rm rev}}\).
Together with the original diagonal normalization, this says
that the marked difference has diagonal \((\zeta,-\zeta)\)
or \((\zeta,0,-\zeta)\).  This is (ALC.13), and the same
formulas prove the actual \(B\)-span assertion directly.
\end{proof}
% END ALC13

% BEGIN ALC14
\paragraph{ALC14. Coordinate-stable subspaces and the rank-three connectivity cases.}
\label{claim:ALC14}
Let \(a:G\to\operatorname{Mat}_n(K)\) be a representation
over a field.  Suppose every \(a(G)\)-stable subspace is
stable under each \(E_{ii}\).  Put an arrow \(j\to i\)
when \(a(g)_{ij}\ne0\) for some \(g\).  If this directed
graph is strongly connected, the representation is irreducible.
For rank two, nonzero entries in both triangular directions
give this connectivity.  For rank three, it follows from
both triangular directions, at least one odd-weight entry,
and the reflection symmetry
\[
 \bigl(\exists g:\ a(g)_{ij}\ne0\bigr)
 \Longleftrightarrow
 \bigl(\exists g:\ a(g)_{j^{\rm rev},i^{\rm rev}}\ne0\bigr).
                                                        \tag{ALC.16}
\]

\begin{proof}
Let \(W\ne0\) be stable.  A nonzero vector in \(W\) has
a nonzero coordinate \(v_j\); applying \(E_{jj}\) and
dividing by \(v_j\) puts \(e_j\) in \(W\).  If \(e_j\)
belongs to \(W\) and \(a(g)_{ij}\ne0\), applying
\(E_{ii}a(g)\) and dividing by that entry puts \(e_i\)
in \(W\).  Strong connectivity therefore puts every basis
vector in \(W\), proving \(W=K^n\).

In rank two the two nonzero triangular directions are exactly
the arrows \(0\to1\) and \(1\to0\).
In rank three, reflection pairs the adjacent lower arrows
\(0\to1\), \(1\to2\), and pairs the adjacent upper
arrows \(2\to1\), \(1\to0\).  Call the existence of
the lower pair \(L\), and of the upper pair \(U\).
The odd-weight condition is exactly \(L\) or \(U\).
If both hold, both directions along the adjacent chain exist.
If \(L\) holds but \(U\) does not, the nonzero upper
triangular direction must be the long arrow \(2\to0\).
Together these form the directed cycle
\(0\to1\to2\to0\).  If \(U\) holds but \(L\) does
not, the lower direction must be \(0\to2\), giving the
opposite directed cycle.  Each case is strongly connected.
These cases exhaust the original hypotheses.
\end{proof}
% END ALC14

% BEGIN ALC15
\paragraph{ALC15. Full original classification of the field-valued types.}
\label{claim:ALC15}
For every field \(K\) and every ring map \(f:A\to K\),
\[
 \operatorname{MatIrr}\bigl(u\otimes f\bigr)
       \Longleftrightarrow \mathrm{TI}(\ker f).
                                                        \tag{ALC.17}
\]
Moreover, if \(\mathfrak m_B\subseteq\pi(\ker f)\),
one of the two triangular entry ideals lies in \(\ker f\).
These are both clauses of the full original
\(\mathrm{types}\) theorem, for ranks two and three and
every original field-valued point.

\begin{proof}
The second assertion is ALC12.  For (ALC.17), an entry ideal
is contained in \(\ker f\) precisely when every one of
its generating entries vanishes after applying \(f\).
If the lower entries all vanish, the first coordinate line
is stable; if the upper entries all vanish, the last
coordinate line is stable.  These are nonzero proper lines
because \(n\ge2\).  In rank three, if the odd-weight
entries all vanish, the middle coordinate line is stable.
Indeed the two entries in its column outside the middle
position have odd weight.  Irreducibility therefore implies
all conditions in (ALC.2).

Conversely suppose (ALC.2) holds.  By ALC13 each coordinate
projector belongs to the actual \(B\)-span of \(u(G_0)\).
After applying \(f\), a subspace stable under the represented
group is stable under those projectors.  Polarization gives
the reflection symmetry (ALC.16): the entry at the reflected
position is the original entry evaluated at \(g^*\).
Applying the same identity twice gives the converse vanishing
implication as well.  The noncontainments in (ALC.2) say
exactly that the graph has the arrows required by ALC14.
That claim proves irreducibility.

No topology, cardinality, or prescribed \(k\)-algebra
structure was imposed on \(K\).  The units used to form
the projectors were already units in \(B\), and their
images remain units under the arbitrary ring map \(f\).
Thus the argument retains the original broad point scope.
It also applies without change when \(\ell=3\).
\end{proof}
% END ALC15

% BEGIN ALC16
\paragraph{ALC16. The continuity, residue and entry-ideal dictionaries.}
\label{claim:ALC16}
For a supplied original local object \(L\), the representation
\(\rho_L\) is adically continuous for \(\mathfrak m_T\).
For every field map \(h:k\to\operatorname{Res}(T)\)
satisfying the original \(\mathrm{IsResMap}(O,T,h)\), its
reduction has the exact \(\mathrm{IsBar}(h)\) diagonal and
ordinary characters.  For every ideal \(P\subset T\),
\[
 \begin{aligned}
 \rho_L.\mathfrak t_+\subseteq P
       &\Longleftrightarrow \mathfrak t_+\subseteq\iota^{-1}P,\\
 \rho_L.\mathfrak t_-\subseteq P
       &\Longleftrightarrow \mathfrak t_-\subseteq\iota^{-1}P,\\
 \rho_L.\mathfrak e\subseteq P
       &\Longleftrightarrow \mathfrak e\subseteq\iota^{-1}P.
 \end{aligned}                                             \tag{ALC.18}
\]
No primality hypothesis on \(P\) is needed.

\begin{proof}
ALC11 gives \(\mathfrak m_B\subseteq\pi(x)\), and the
original \(L.\mathrm{lies}\) gives
\(x=\iota^{-1}\mathfrak m_T\).  Therefore
\(\iota(x_0^N)\subseteq\mathfrak m_T^N\) for every
\(N\).  Compose each locally constant quotient map in
RTC9 with the induced quotient map to
\(T/\mathfrak m_T^N\).  This proves continuity of every
matrix entry and every ordinary-character value, with all
original local group arguments and all powers \(N\).

For the exact residue assertion, let \(b\in B\), and
choose \(a\in O\) lifting its residue in the original
\(\mathcal R\).  Then \(b-a\in\mathfrak m_B\), so
\(\iota(b-a)\in\mathfrak m_T\).  The prescribed
compatibility of \(h\) with \(O\) gives
\[
 \operatorname{res}_T(\iota(b))
       =h\bigl(\operatorname{res}_{\mathcal R}(b)\bigr).
                                                        \tag{ALC.19}
\]
Every diagonal entry and every ordinary-character value is
an actual element of \(B\).  Their residues in the supplied
space are, respectively, \(\Psi(g)^{w(n)_i}\) and
\(\Psi(\operatorname{loc}_v s)^{w(n)_i}\).
Applying (ALC.19) is exactly the two-part
\(\mathrm{IsBar}(h)\) assertion.  No equality of
semisimplifications is substituted for these literal entries.

Finally, an ideal generated by the \(\iota\)-images of
a family lies in \(P\) if and only if every original
generator lies in \(\iota^{-1}P\).  Apply this statement
to each of the three defining entry predicates.  It gives
(ALC.18) for every ideal, including the zero and unit ideals.
\end{proof}
% END ALC16

% BEGIN ALC17
\paragraph{ALC17. Full original local representation theorem.}
\label{claim:ALC17}
For every original \(L:S.\mathrm{Loc}(x)\), all clauses
of \(\mathrm{loc\_rho}\) hold.  In addition to the
continuity, every specified residue map, and all ideal
comparisons in ALC16, its final assertion is
\[
 \mathfrak t_+\not\subseteq x
 \quad\Longrightarrow\quad
 \rho_L.\mathfrak t_+=T,
 \qquad \rho_L.\mathfrak t_-\subseteq\mathfrak m_T.
                                                        \tag{ALC.20}
\]
The statement holds for both allowed ranks on the same
original \(F,S,L\).

\begin{proof}
All preceding fields are ALC16.  If
\(\rho_L.\mathfrak t_+\) were a proper ideal of the local
ring \(T\), it would lie in \(\mathfrak m_T\).
The first equivalence in (ALC.18) and
\(\iota^{-1}\mathfrak m_T=x\) would give
\(\mathfrak t_+\subseteq x\), a contradiction.
Thus \(\rho_L.\mathfrak t_+=T\).

Consider the actual field-valued map
\(A\xrightarrow{\iota}T\to\operatorname{Res}(T)\).
Its kernel is \(x\), and ALC11 gives the required
\(\mathfrak m_B\) contraction.  The special-fiber clause
of the full types theorem, ALC15, says that one triangular
entry ideal is contained in \(x\).  The hypothesis excludes
the lower one, so \(\mathfrak t_-\subseteq x\).
The second equivalence in (ALC.18) gives the remaining
assertion in (ALC.20).  In rank three this uses the proved
general opposite-weight product argument, not a rank-two
point-realization shortcut.
\end{proof}
% END ALC17

\subsection{Existence at every original closed point}

% BEGIN ALC18
\paragraph{ALC18. The original closed point has a finite residue field.}
\label{claim:ALC18}
If \(x\subset A\) is maximal and
\(\mathfrak m_B\subseteq\pi(x)\), then \(K=A/x\)
is a finite field with its prescribed \(O\)-algebra
structure, and \(\mathfrak m_O\) acts by zero on it.

\begin{proof}
RTC7 makes \(A\) a finite-type algebra over the local
coefficient ring \(B\), whose residue field is the original
finite \(k\).  The contraction hypothesis makes \(K\)
a finite-type \(k\)-algebra, and maximality makes it a
field.  We recall the finite-type field argument to make
the residue extension explicit.

For a field generated as a \(k\)-algebra by finitely many
elements, choose a maximal algebraically independent subset
\(t_1,\ldots,t_d\) of these generators.  The remaining
generators are algebraic over \(k(t_1,\ldots,t_d)\).
Clearing the finitely many denominators of their monic
equations makes them integral over
\(k[t_1,\ldots,t_d,h^{-1}]\) for one nonzero polynomial
\(h\).  The whole generated field is a finite integral
algebra over this ring, since it already contains \(h^{-1}\).
A subring over which a field is integral is itself a field:
the monic equation for the inverse of a nonzero element
expresses that inverse in the subring.

If \(d>0\), choose an irreducible polynomial in \(k[t_1]\)
which does not divide \(h\).  Such a polynomial exists
because a one-variable polynomial ring over a field has
infinitely many irreducibles, by the usual product-plus-one
argument.  It remains prime after adjoining the other
variables and is not made invertible by inverting \(h\).
Thus \(k[t_1,\ldots,t_d,h^{-1}]\) is not a field, a
contradiction.  Hence \(d=0\), so the original field is
a finite algebraic extension of \(k\), and is finite as
a set because \(k\) is finite.

Apply this to \(K=A/x\).  The local map \(O\to B\)
sends \(\mathfrak m_O\) into \(\mathfrak m_B\), whose
image in \(K\) is zero.  The resulting \(k\)-embedding
is the one induced by the original \(O\)-structure; no
replacement of that embedding has been made.
\end{proof}
% END ALC18

% BEGIN ALC19
\paragraph{ALC19. Constructing the required finite unramified coefficient extension.}
\label{claim:ALC19}
Let \(O\) be any original coefficient DVR and let \(K\)
be a finite field with an \(O\)-algebra structure killing
\(\mathfrak m_O\).  There is an explicitly constructed
finite free \(O\)-algebra \(O'\), itself a coefficient
DVR of residue characteristic \(\ell\), such that
\[
 \mathfrak m_OO'=\mathfrak m_{O'},\qquad
              K\xrightarrow{\ \sim\ }\operatorname{Res}(O')
                  \quad\hbox{as \(O\)-algebras}.
                                                        \tag{ALC.21}
\]
This supplies the full original \(\mathrm{unr\_ex}\)
output, with an isomorphism in place of its required embedding.

\begin{proof}
The given map factors through an embedding of the finite
field \(k=\operatorname{Res}(O)\).  Choose a generator
\(\gamma\) of the cyclic group \(K^\times\); it also
generates \(K\) as a field over \(k\).  For completeness,
cyclicity follows from the elementary finite-subgroup argument:
in a finite abelian subgroup of a field's units, an element
has order equal to the group exponent, by combining its
prime-primary maximal orders.  The polynomial \(X^e-1\)
has at most \(e\) roots, so this exponent is at least the
group order and therefore equal to it.

Let \(g\in k[X]\) be the monic minimal polynomial of
\(\gamma\).  Lift its coefficients to \(O\), retaining
the leading coefficient one, to obtain a monic polynomial
\(G\) of the same positive degree.  Define
\[
                    O'=O[X]/(G).
                                                        \tag{ALC.22}
\]
Division by the monic polynomial gives the finite free
\(O\)-basis \(1,X,\ldots,X^{\deg G-1}\).  Let
\(\varpi\) be a uniformizer of \(O\).  This basis shows
that \(O'\) is complete and separated for \(\varpi O'\),
and that multiplication by \(\varpi\) is injective.  Its
quotient is exactly
\[
          O'/\varpi O'=k[X]/(g)\simeq K.
                                                        \tag{ALC.23}
\]
An element outside \(\varpi O'\) has an inverse modulo
that ideal.  The convergent geometric series corrects it
to an actual inverse.  Thus \(O'\) is local with maximal
ideal \(\varpi O'\).

Every nonzero element \(b\in O'\) has a finite maximal
power \(\varpi^r\) dividing it: in the displayed free
basis, take the minimum valuation of its nonzero
coefficients.  Thus \(b=\varpi^r u\) with \(u\) a
unit.  Since \(\varpi\) is a non-zero-divisor, products
of nonzero elements are nonzero.  The ring is therefore
a domain.  In any nonzero ideal, choose an element of
smallest such exponent \(r\); the ideal is then exactly
\(\varpi^rO'\).  This proves the discrete valuation
property and Noetherianity directly.  The uniformizer is
nonzero and not a unit, so the ring is not a field.

Its maximal-ideal completeness is the already proved
\(\varpi\)-adic completeness.  Its residue is finite and
has characteristic \(\ell\), proving all coefficient
clauses.  The image of \(\mathfrak m_O=(\varpi)\) is
exactly the maximal ideal.  Finally the inverse of the
specific isomorphism in (ALC.23), sending the polynomial
class to \(\gamma\), is the map in (ALC.21).  The
coefficient lift and this map respect the prescribed
\(O\)-action throughout.
\end{proof}
% END ALC19

% BEGIN ALC20
\paragraph{ALC20. The residue-compatible coefficient completion and its universal property.}
\label{claim:ALC20}
Let \(O\) be any commutative local ring, \(D\) a
Noetherian \(O\)-algebra, and
\(\varphi:D\to\operatorname{Res}(O)\) an
\(O\)-algebra map.  There are one coefficient algebra
\(R\) and one map \(j:D\to R\) such that
\[
          \operatorname{res}_R(j(d))=\varphi(d),
                                                        \tag{ALC.24}
\]
and every residue-compatible map from \(D\) to a coefficient
algebra factors uniquely through \(j\).  The same conclusion
holds for any already supplied \(\mathrm{IsCplAt}(\ker\varphi,R)\)
with its given \(O\)-structure and scalar tower.  These
are the full original \(\mathrm{cpl\_ex}\) and
\(\mathrm{cpl\_univ}\) statements.

Independently of a chosen coefficient ring \(O\), let
\(D\) be any Noetherian ring, \(y\) any maximal ideal,
and \(R\) any local \(D\)-algebra satisfying
\(\mathrm{IsCplAt}(y,R)\).  For every commutative \(T\),
every ideal \(J\) for which \(T\) is complete and
separated, and every \(f:D\to T\) with \(f(y)\subseteq J\),
there is a unique extension sending \(\mathfrak m_R\)
into \(J\); moreover
\(R\simeq_D\mathrm{Cpl}(y)\) and \(yR=\mathfrak m_R\).
These are all outputs of the full original
\(\mathrm{cpl\_uniq}\).

\begin{proof}
The map \(\varphi\) is surjective on the image of \(O\),
so \(y=\ker\varphi\) is maximal.  Construct
\(R=\varprojlim_N D/y^N\) and the canonical \(j\).
ALC5 proves that \(R\) is Noetherian local and complete,
with maximal ideal \(yR\) and residue \(D/y\).
The map induced by \(\varphi\) identifies that residue
with the original \(k\).  The composite \(O\to R\)
is local and surjective on residues: a nonunit has residue
zero, and every residue is lifted from \(O\).  Thus it
is the required coefficient algebra and (ALC.24) holds.

If \(f:D\to C\) has the stated residue, then
\(f(y^N)\subseteq\mathfrak m_C^N\) for every \(N\).
The compatible maps
\(D/y^N\to C/\mathfrak m_C^N\), followed by completeness
of \(C\), define an \(O\)-algebra map \(g:R\to C\).
It extends \(f\).  Any other such \(O\)-algebra map is
local and continuous by ALC1.  Since \(j(D)\) is dense,
the two maps agree modulo every \(\mathfrak m_C^N\),
and separatedness gives equality.  The objects \(R,j\)
were fixed before \(C,f\), as the original contract requires.

For an already supplied \(R\) with
\(\mathrm{IsCplAt}(y,R)\), its exact contraction and
density give \(R/\mathfrak m_R\simeq D/y\), respecting
the given scalar tower.  This makes the structure map
\(O\to R\) local with the required residue field, so
\(R\) is a coefficient algebra.  Its compatible finite
quotients identify it, with the original map from \(D\),
with the same inverse limit.  The preceding construction
and uniqueness therefore apply to that ring and that map,
proving the wider \(\mathrm{cpl\_univ}\) assertion too.

For the coefficient-independent assertion, the same exact
quotient and density argument identifies the given \(R\)
with \(\varprojlim D/y^N\).  No local-ring or coefficient
structure on \(T\) is needed.  The maps
\(D/y^N\to T/J^N\) define the extension by \(J\)-adic
completeness.  The identity \(yR=\mathfrak m_R\) follows
from this inverse-limit identification and (ALC.6); it
shows that the extension sends \(\mathfrak m_R\) into
\(J\).  Any two extensions with that property agree
modulo \(J^N\), by approximating each \(r\in R\) with
an element of \(D\) modulo \(\mathfrak m_R^N\).
Separatedness gives uniqueness.  Finally ALC5 identifies
this inverse limit with the actual \(\mathrm{Cpl}(y)\),
as a \(D\)-algebra.  This proves the claimed full
\(\mathrm{cpl\_uniq}\) scope.
\end{proof}
% END ALC20

% BEGIN ALC21
\paragraph{ALC21. Full original existence of the arithmetic local object.}
\label{claim:ALC21}
For every original maximal ideal \(x\subset A\) with
\(\mathfrak m_B\subseteq\pi(x)\), an original
\(L:S.\mathrm{Loc}(x)\) exists.  This is the full
\(\mathrm{loc\_ex}\) theorem at both allowed ranks.

\begin{proof}
By ALC18, \(K=A/x\) is a finite field with the required
\(O\)-structure.  ALC19 constructs a finite coefficient
DVR extension \(O'\), with
\(\mathfrak m_OO'=\mathfrak m_{O'}\), and an
\(O\)-algebra embedding \(e:K\to k'\), where
\(k'=\operatorname{Res}(O')\).  Keep this choice fixed.
Put \(D=O'\otimes_O A\), a Noetherian ring because it
is finite over the Noetherian ring \(A\).  Define
\[
 \varphi:D\longrightarrow k',\qquad
       b\otimes a\longmapsto\operatorname{res}_{O'}(b)
                                  e(a\bmod x).
                                                        \tag{ALC.25}
\]
The prescribed scalar compatibility makes this a well-defined
\(O'\)-algebra map.  ALC20 constructs its coefficient
completion \(j:D\to T\).  Give \(T\) the composite
\(O\)-structure and put
\[
                     \iota(a)=j(1\otimes a).
                                                        \tag{ALC.26}
\]
Then
\(\operatorname{res}_T(\iota(a))=e(a\bmod x)\).
Injectivity of \(e\) gives
\(\iota^{-1}\mathfrak m_T=x\), the literal
\(\mathrm{lies}\) field.

It remains to prove the exact \(\mathrm{univ}\) field,
including its scalar structures.  Let \(C\) be any
coefficient algebra over \(O'\), with the supplied
\(O\)-structure and scalar tower, and let \(f:A\to C\)
be an \(O\)-algebra map whose residue agrees with
\(\iota\).  The tensor map
\[
               f':D\to C,\qquad b\otimes a\mapsto b f(a)
\]
is \(O'\)-linear.  Its residue is (ALC.25), first on
pure tensors by the assumed equality for \(a\), then
on their sums.  The universal property in ALC20 supplies
a unique \(O'\)-algebra map \(g:T\to C\) with
\(gj=f'\).  Restriction to \(1\otimes A\) gives
\(g\iota=f\).

If another \(O'\)-algebra map has this equality on \(A\),
then its composite with \(j\) agrees with \(f'\) on
every pure tensor, and hence on \(D\).  Uniqueness in
ALC20 forces it to equal \(g\).  This proves exactly the
original local universal property.  The finite coefficient
extension, its unramified maximal-ideal identity, all
coefficient-algebra fields of \(T\), both scalar structures
and their tower, \(\iota\), \(\mathrm{lies}\), and
\(\mathrm{univ}\) have now all been constructed on the
same object.
\end{proof}
% END ALC21

% BEGIN ALC22
\paragraph{ALC22. The actual local component charts are now available to propagation.}
\label{claim:ALC22}
For every original closed point \(x\) over
\(\mathfrak m_B\), the local object constructed in ALC21
has all the completion, flatness and representation properties
proved above.  The actual spectral map
\[
       f_L:\operatorname{Spec}T\longrightarrow\operatorname{Spec}A,
                     \qquad P\longmapsto\iota^{-1}P
                                                        \tag{ALC.27}
\]
is the full component chart required by the original
\(\mathrm{ChartOfFlat}\) consumer: it is continuous,
sends the closed point to \(x\), sends minimal primes to
generic points of irreducible components, and realizes every
irreducible component through \(x\) in this way.  For a
core point, its local lower-entry ideal is the unit ideal
and its upper-entry ideal lies in \(\mathfrak m_T\).

\begin{proof}
The construction supplies an actual \(L\), so ALC4,
ALC11 and ALC17 apply without another existence premise.
In particular \(A\to T\) is flat and
\(\iota^{-1}\mathfrak m_T=x\).  Apply the already
proved ACP16 with the identity open chart of
\(\operatorname{Spec}A\).  Its assumptions are exactly
these properties, and its map is (ALC.27), with no change
of the underlying algebra structure.  This gives every
component-chart clause.  Alternatively the prime and
minimal-prime argument in ALC9 is the same going-down
mechanism after the coefficient tensor extension.

For a point of the original core, the extra condition is
\(\mathfrak t_+\not\subseteq x\).  ALC17 gives the
two asserted local triangular ideals.  The ideal comparisons
of ALC16 identify the local \(\mathrm{TI}\), exceptional
entry and good loci with the inverse images of their
original global conditions.  In particular, for an original
ideal \(M\subset B\),
\[
          M\subseteq L.\pi(P)
             \Longleftrightarrow
          M\subseteq\pi(\iota^{-1}P),
                                                        \tag{ALC.28}
\]
and the right side is the original global good condition.
Thus the charts fit the actual local-to-global comparison
used in \(\mathrm{ComponentsGoodRankTwo}\), and its
rank-three counterpart, with the same rings and ideals.

This does not supply the remaining numerical and arithmetic
fields of their propagation records.  The sharp component
chains, dimension bounds, nice-prime hypotheses, global
density and seed incidence must still be supplied.  The
conclusion here is the completed original local producer,
not the whole propagation theorem.
\end{proof}
% END ALC22

\paragraph{Scope of the completed producer.}
The full original \(\mathrm{cpl\_of\_univ}\),
\(\mathrm{loc\_cpl}\), \(\mathrm{c3a}\),
\(\mathrm{cpl\_is}\), \(\mathrm{cpl\_uniq}\),
\(\mathrm{loc\_flat}\), \(\mathrm{idem}\),
\(\mathrm{types}\), \(\mathrm{loc\_rho}\),
\(\mathrm{unr\_ex}\), \(\mathrm{cpl\_univ}\),
\(\mathrm{cpl\_ex}\), and \(\mathrm{loc\_ex}\)
are supplied at the scopes stated above.  The elementary
completion lemmas allow arbitrary Noetherian base rings;
the coefficient-completion universal lemmas allow arbitrary
commutative local coefficient rings.  The arithmetic
theorems retain the original DVR, \(F.\mathrm{Ok}\),
rank-two-or-three, and supplied \(\mathrm{Space}\) premises.
No finite-point conclusion has been used upstream.

Rank-two and rank-three spaces and their coefficient ideals
remain separate.  Their local coefficient extensions need
not be identified.  All constructions keep the original
frame, group maps, ordinary characters and marked element;
the original arithmetic \(P\) premise remains where the
AFE entry uses it.  The allowed Host case \(\ell=3\) is
preserved.  Existence of the supplied universal spaces,
sharp dimension and component bounds, closed-point and
seed incidence, the rank-three corner and exceptional
ideal, full propagation, \(\mathrm{OffFin}\), automorphic
seeds and patching/classicality remain separate.  The
stronger Alt outputs still retain their original
\(\mathrm{WitA}\) and \(\mathrm{WitAQ}\) producers.
No whole selected RH/GRH obligation is closed by this chapter.

\subsection{Exact source correspondence}

All primary files are at the accepted commit
\texttt{7cbf7a0c32fca1badf05931879a37c3f8a2b03b4}.
Write \(F\) for
\texttt{src/\allowbreak{}Development/\allowbreak{}Automorphic\allowbreak{}Arithmetic/\allowbreak{}}
\texttt{Automorphic\allowbreak{}Varieties/\allowbreak{}Frobenius\allowbreak{}Smooth\allowbreak{}Vectors/\allowbreak{}}
and \(H\) for the sibling
\texttt{Commutative\allowbreak{}Algebra\allowbreak{}Height/\allowbreak{}} directory.
Full snapshots, exact ranges and hashes are in this component's
source manifest.

\begin{itemize}
\item ALC1--ALC4 match \(F\)\texttt{LocCpl.lean},
lines 7--188, especially the complete universal theorem at
100--161 and actual arithmetic application at 162--188.
The original \(\mathrm{IsCO}\) record is in the sibling
\texttt{CommutativeAlgebra/ResMapBijective.lean}, lines 9--13;
\(\mathrm{IsCplAt}\) is in \(H\)\texttt{IsLocalHom.lean}.
\item ALC5--ALC9 reconstruct the completion and flatness
mathematics used by \(H\)\texttt{CplIs.lean},
\texttt{CplUniq.lean}, and
\texttt{HighestWeightSpaceFinite.lean}; the last contains
the full \(\mathrm{c3a}\) statement and its prime clauses.
\item ALC10--ALC11 give all of \(F\)\texttt{LocFlat.lean},
lines 9--69.  ALC12--ALC15 give the full \(\mathrm{types}\)
statement and its supporting matrix arguments in
\(F\)\texttt{IrrOfConn.lean}, with the original projector
theorem in \texttt{AddMulSelfEqOfOrthogonalIdempotents.lean}.
\item ALC16--ALC17 give all outputs of
\(F\)\texttt{LocRho.lean}, lines 8--66, including every
ideal, every specified residue map, and both ranks in its
last implication.
\item ALC18--ALC21 reconstruct the source suppliers and
full existence theorem in \(F\)\texttt{ResEqOf.lean},
lines 11--178.  The monic-polynomial construction in ALC19
proves the needed unramified coefficient extension directly.
\item ALC22 applies the unchanged ACP16 to the actual map
used in the sibling
\texttt{CurveDVRSpecialization/\allowbreak{}ComponentsGoodRankTwo.lean},
lines 107--128 and 149--158.  The remaining dimension,
density and seed inputs in that source are retained explicitly.
\end{itemize}
